\documentclass[UTF8,a4paper,10pt]{ctexart}
\usepackage[a4paper,top=14mm,bottom=16mm,left=16mm,right=16mm,footskip=8mm]{geometry}
\usepackage{amsmath,amssymb,booktabs,tabularx,array,enumitem,fancyhdr,lastpage,xcolor,hyperref,multirow}
\setCJKmainfont{SimSun}\setCJKsansfont{SimHei}\setmainfont{Times New Roman}
\hypersetup{hidelinks,unicode=true}\pagestyle{fancy}\fancyhf{}\fancyfoot[C]{第\thepage 页\quad 共\pageref*{LastPage}页}\renewcommand{\headrulewidth}{0pt}
\renewcommand{\arraystretch}{1.2}\setlist[enumerate]{leftmargin=2.2em,itemsep=.25em,topsep=.2em}
\newcommand{\sect}[1]{\par\smallskip\noindent\fcolorbox{black}{black!4}{\parbox{\dimexpr\linewidth-2\fboxsep-2\fboxrule}{\sffamily\bfseries #1}}\par\smallskip}
\newcommand{\opts}[1]{\par\hspace*{.5em}\begin{tabularx}{\dimexpr\linewidth-.5em}{@{}XXXX@{}}#1\end{tabularx}\par}
\begin{document}\small
\begin{center}{\zihao{2}\sffamily\bfseries 湖南工商大学课程考核试卷【问卷】（B卷）}\par{\zihao{3}\bfseries 统计学B补考卷}\end{center}
\noindent 考生务必将答案写在答卷上，问卷上答题一律无效；考生可用普通计算器，结果保留4位小数。
\sect{一、单选题（每小题1分，共10分）}
\begin{enumerate}
\item 构成统计总体的最原始的承担者是（）。\opts{A.总体&B.总体单位&C.调查单位&D.法人单位}
\item 研究某工业企业职工的工资水平时，该企业职工的工资总额是（）。\opts{A.数量标志&B.品质标志&C.质量指标&D.数量指标}
\item 对于生产过程中产品质量的检查和控制宜采用（）。\opts{A.普查&B.重点调查&C.典型调查&D.抽样调查}
\item 下列属于总量指标的是（）。\opts{A.职工数&B.及格率&C.达标率&D.出勤率}
\item 某国某年的钢产量与人口总数之比是（）。\opts{A.比较相对数&B.比例相对数&C.强度相对数&D.平均数}
\item 反映总体分布离中趋势的是（）。\opts{A.变异指标&B.相对指标&C.平均指标&D.总量指标}
\item 定基发展速度等于相应时期各个环比发展速度的（）。\opts{A.和&B.差&C.积&D.商}
\item 以年度为周期的规律性变动是（）。\opts{A.长期趋势&B.季节变动&C.循环变动&D.不规则变动}
\item \texttt{TODO：原照片第9题题干严重失焦，仅能辨认选项为拉氏、帕氏、杨氏、费氏，题干不作猜写。}
\item 下列抽样方法中，抽样误差最小的是（其他条件相同）（）。\opts{A.纯随机抽样&B.等距抽样&C.分层抽样&D.整群抽样}
\end{enumerate}
\sect{二、判断题（每小题1分，共10分）}
\begin{enumerate}[start=11]
\item 产品产量可以是连续变量，也可以是离散变量。（）
\item 重点调查中的重点单位是指这些单位是工作中的重点。（）
\item 对于连续型变量一般只能编制组距数列。（）
\item 不管是按品质标志分组，还是按数量标志分组，都能形成分配数列。（）
\item 相对指标可以使不能直接对比的总量指标找到对比的基础。（）
\item 男女性别比属于结构相对数。（）
\item 总体中各标志值之间的差异程度越大，标准差系数就越小。（）
\item 若一个指数的同度量因素是数量指标，则这个指数就是数量指标指数。（）
\item 只有当相关系数接近$+1$时，才能说明两变量之间存在高度相关关系。（）
\item 对于两个变量而言，其相关系数、协方差和回归系数的符号总是一致。（）
\end{enumerate}
\sect{三、填空题（每小题4分，共20分）}
\begin{enumerate}[start=21]
\item 某公司生产成本计划比上年降低10\%，实际降低20\%。则该公司生产成本的计划完成程度为\underline{\hspace{2.5cm}}，是否完成了任务（选填是或否）\underline{\hspace{1.5cm}}。
\item 适度偏态下某数列的算术平均数$\bar x=31$，中位数$M_e=29$。则该数列的众数$M_o$的大约值为\underline{\hspace{2cm}}，该数列的分布特征为（选填左偏或右偏）\underline{\hspace{2cm}}。
\item 已知$n=10,\sum x=71,\sum y=288,\sum x^2=601,\sum y^2=11612,\sum xy=2454$，则$x,y$的相关系数$r$为\underline{\hspace{2cm}}，相关程度为（选填低度、显著或高度相关）\underline{\hspace{2cm}}。
\item 某数列的方差为47，各变量值平方的平均数为672。则该数列的算术平均数为\underline{\hspace{2cm}}，标准差系数为\underline{\hspace{2cm}}。
\item 某企业2010--2017年间产量年均增长率为14\%，2017--2020年间产量总发展速度为182\%，2024年产量是2020年的2倍。则2010--2024年间该企业产量年均发展速度为\underline{\hspace{2cm}}，年均增长率为\underline{\hspace{2cm}}。
\end{enumerate}
\sect{四、计算题（每小题12分，共36分）}
\begin{enumerate}[start=26]
\item \texttt{TODO：原文件没有完整拍到第26题题面；第3页顶部仅残留该题末尾三项要求，无法可靠恢复题干和数据。}
\item 某企业商品销售额如下表（单位：万元），请完成下列计算。
\begin{center}\begin{tabular}{crrrrr}\toprule 年份&第一季度&第二季度&第三季度&第四季度&总平均或合计\\\midrule2023年&280&560&140&190&1170\\2024年&310&455&190&160&1115\\2025年&295&510&210&180&1195\\平均&&&&&\\季节指数&&&&&\\2026年&&&&&1200\\\bottomrule\end{tabular}\end{center}
\begin{enumerate}[label=（\arabic*）,leftmargin=2.8em]\item 计算各季度平均值、三年的季度总平均值，填入表格；（4分）\item 计算季节指数、季节指数之和，填入表格；（4分）\item 2026年计划全年销售额1200万元，计算各季应完成销售额，填入表格。（4分）\end{enumerate}
\item 某商场甲、乙、丙三种食品2025年一季度和二季度的销售额和价格资料如下表（单位：元），请完成下列计算。
\begin{center}\begin{tabular}{crrrrrr}\toprule\multirow{2}{*}{品种}&\multicolumn{2}{c}{销售额}&\multicolumn{2}{c}{价格}&\multicolumn{2}{c}{计算栏}\\\cmidrule(lr){2-3}\cmidrule(lr){4-5}\cmidrule(lr){6-7}&$p_0q_0$&$p_1q_1$&$p_0$&$p_1$&$K_q$&$p_1q_1/K_q$\\\midrule 甲&1800&2100&10&12&&\\乙&2700&2500&20&25&&\\丙&3600&4900&18&24&&\\合计&8100&9500&&&&\\\bottomrule\end{tabular}\end{center}
\begin{enumerate}[label=（\arabic*）,leftmargin=2.8em]\item 计算相关数据，填入表格计算栏；（3分）\item 计算三种商品的销售额总指数、价格总指数、销售量总指数；（6分）\item 计算价格和销量变动对销售额变动的影响额。（3分）\end{enumerate}
\end{enumerate}
\sect{五、综合题（每小题12分，共24分）}
\begin{enumerate}[start=29]
\item 某班50名学生统计学成绩如下表（单位：分）。
\begin{center}\begin{tabular}{*{10}{r}}39&47&58&61&62&63&65&64&64&72\\64&65&68&67&69&69&70&70&71&73\\72&73&74&75&75&75&76&76&76&77\\77&78&78&79&79&79&79&79&82&83\\84&85&85&86&86&87&88&88&95&98\end{tabular}\end{center}
\begin{enumerate}[label=（\arabic*）,leftmargin=2.8em]\item 根据原始数据编制组距式变量数列，填入如下表格：（4分）
\begin{center}\begin{tabular}{ccccc}\toprule 成绩&组中值$x$&次数$f$&频率&$xf$\\\midrule60以下&&&&\\60--70&&&&\\70--80&&&&\\80--90&&&&\\90以上&&&&\\合计&&&&\\\bottomrule\end{tabular}\end{center}
\item 根据变量数列计算成绩的加权算术平均数、中位数和众数；（6分）\item 根据算术平均数、中位数和众数分析该数列的分布特征。（2分）\end{enumerate}
\item 我国居民2005--2024年人均国内生产总值$x$与农村居民人均可支配收入$y$数据（单位：百元）的Excel回归分析结果如下表。
\begin{center}\begin{tabular}{lr}\multicolumn{2}{c}{回归统计}\\\toprule Multiple R&0.9977\\R Square&0.9953\\Adjusted R Square&0.9951\\标准误差&4.4188\\观测值&20\\\bottomrule\end{tabular}\par\smallskip
\begin{tabular}{lrrrrr}\multicolumn{6}{c}{方差分析}\\\toprule &df&SS&MS&F&Significance F\\回归分析&1&74720.48&74720.48&3826.74&$2.00835\mathrm E{-22}$\\残差&18&351.47&B&&\\总计&A&75071.94&&&\\\bottomrule\end{tabular}\par\smallskip
\begin{tabular}{lrrrr}\toprule &Coefficients&标准误差&t Stat&P-value\\Intercept&$-9.7093$&2.2647&$-4.2873$&0.000443448\\X Variable 1&0.2424&0.0039&61.8607&$2.00835\mathrm E{-22}$\\\bottomrule\end{tabular}\end{center}
\begin{enumerate}[label=（\arabic*）,leftmargin=2.8em]\item 写出相关系数和估计标准误差；（2分）\item 写出表格中的阴影部分（A、B）的数值；（4分）\item 写出人均国内生产总值与农村居民人均可支配收入的一元线性回归方程；（4分）\item 若2025年人均国内生产总值1100百元，预测该年农村居民人均可支配收入。（2分）\end{enumerate}
\end{enumerate}
\newpage\sect{校对说明}
原PDF未附参考答案。第26题以及第9题题干在源文件中本身缺失或失焦，已明确标记TODO；其余大题的原始数据、计算栏、分组表、回归表和全部分问均按可见原页重建，没有用概括句替代。
\end{document}
